\documentclass[12pt,a4paper]{article}

% --------------------------------------------------
% Язык, шрифты и кодировка для PDFLaTeX
% --------------------------------------------------
\usepackage{cmap}
\usepackage[utf8]{inputenc}
\usepackage[T2A]{fontenc}
\usepackage[english,russian]{babel}
\usepackage{lmodern}

% --------------------------------------------------
% Настройки страницы
% --------------------------------------------------
\usepackage[
    left=25mm,
    right=20mm,
    top=22mm,
    bottom=22mm
]{geometry}

\usepackage{setspace}
\onehalfspacing

\usepackage{indentfirst}
\setlength{\parindent}{1.25cm}
\setlength{\parskip}{0.2em}

% --------------------------------------------------
% Пакеты для оформления
% --------------------------------------------------
\usepackage{graphicx}
\usepackage{float}
\usepackage{booktabs}
\usepackage{array}
\usepackage{longtable}
\usepackage{caption}
\usepackage{subcaption}
\usepackage{xcolor}
\usepackage{hyperref}
\usepackage{enumitem}
\usepackage{fancyhdr}
\usepackage{titlesec}
\usepackage{tcolorbox}
\usepackage{listings}

% --------------------------------------------------
% Цвета
% --------------------------------------------------
\definecolor{mainblue}{HTML}{1F4E79}
\definecolor{lightblue}{HTML}{EAF2F8}
\definecolor{darkgray}{HTML}{333333}
\definecolor{codebg}{HTML}{F5F5F5}

% --------------------------------------------------
% Гиперссылки
% --------------------------------------------------
\hypersetup{
    colorlinks=true,
    linkcolor=mainblue,
    urlcolor=mainblue,
    citecolor=mainblue
}

% --------------------------------------------------
% Заголовки
% --------------------------------------------------
\titleformat{\section}
{\Large\bfseries\color{mainblue}}
{\thesection.}
{0.5em}
{}

\titleformat{\subsection}
{\large\bfseries\color{darkgray}}
{\thesubsection.}
{0.5em}
{}

% --------------------------------------------------
% Колонтитулы
% --------------------------------------------------
\pagestyle{fancy}
\fancyhf{}
\lhead{Food Delivery Analytics}
\rhead{\thepage}
\setlength{\headheight}{15pt}
\renewcommand{\headrulewidth}{0.4pt}

% --------------------------------------------------
% SQL-код
% --------------------------------------------------
\lstdefinestyle{sqlstyle}{
    backgroundcolor=\color{codebg},
    basicstyle=\ttfamily\small,
    breaklines=true,
    frame=single,
    rulecolor=\color{lightblue},
    keywordstyle=\color{mainblue}\bfseries,
    commentstyle=\color{gray},
    stringstyle=\color{darkgray},
    showstringspaces=false
}

\lstset{style=sqlstyle}

% --------------------------------------------------
% Команда для вставки графиков
% Если картинки нет, в PDF появится аккуратная заглушка.
% --------------------------------------------------
\newcommand{\insertfigure}[3]{
    \begin{figure}[H]
        \centering
        \IfFileExists{#1}{
            \includegraphics[width=0.92\textwidth]{#1}
        }{
            \fbox{
                \begin{minipage}[c][5cm][c]{0.85\textwidth}
                    \centering
                    Здесь должна быть картинка.\\
                    Добавьте файл изображения в папку \texttt{images/}.
                \end{minipage}
            }
        }
        \caption{#2}
        \label{#3}
    \end{figure}
}

\begin{document}

% --------------------------------------------------
% Титульная страница
% --------------------------------------------------
\begin{titlepage}
    \centering

    \vspace*{2cm}

    {\Huge\bfseries\color{mainblue} Food Delivery Product \& Courier Analytics \par}

    \vspace{0.8cm}

    {\Large Аналитический проект по данным сервиса доставки еды \par}

    \vspace{1.5cm}

        \centering
        \large
        SQL-анализ продуктовых метрик, эффективности доставки, повторных заказов и удержания пользователей.

    \vfill

    \begin{flushleft}
        \large
        \textbf{Автор:} Затравин Матвей \\
        \textbf{Инструменты:} SQL, SQLite, DBeaver, Python, Matplotlib \\
        \textbf{Тип проекта:} Проект в портфолио для позиции продуктовый аналитик \\
        \textbf{Год:} 2026
    \end{flushleft}

    \vspace{1cm}

\end{titlepage}

\tableofcontents
\newpage

% --------------------------------------------------
\section{Описание проекта}

Проект посвящён анализу данных сервиса доставки еды. В работе рассматриваются две стороны продукта: пользовательская активность и операционная эффективность доставки.

Основная цель проекта -- рассчитать ключевые продуктовые и операционные метрики, найти сильные и слабые места сервиса, а также предложить рекомендации, которые могут помочь улучшить повторные заказы и удержание пользователей.

В анализе использовались таблицы с действиями пользователей и курьеров. На основе этих данных были рассчитаны delivery rate, среднее время доставки, пиковые часы нагрузки, эффективность курьеров, repeat order rate и retention после первого успешного заказа.

% --------------------------------------------------
\section{Цели анализа}

В рамках проекта были поставлены следующие аналитические задачи:

\begin{enumerate}
    \item Оценить долю успешно доставленных заказов.
    \item Рассчитать среднее время доставки.
    \item Определить, в какие часы нагрузка на доставку максимальна.
    \item Проанализировать различия в эффективности курьеров.
    \item Рассчитать долю пользователей с повторными заказами.
    \item Провести retention-анализ после первого успешного заказа.
    \item Проверить, связаны ли пол и возраст курьера со средним временем доставки.
    \item Сформулировать продуктовые и операционные рекомендации на основе полученных данных.
\end{enumerate}

% --------------------------------------------------
\newpage
\section{Данные}

В проекте использовались три основные таблицы:

\begin{itemize}
    \item \texttt{courier\_actions} -- действия курьеров с заказами;
    \item \texttt{user\_actions} -- действия пользователей с заказами;
    \item \texttt{couriers} -- профиль курьеров: дата рождения и пол.
\end{itemize}

Основные поля таблиц:

\begin{table}[H]
\centering
\caption{Описание основных полей}
\begin{tabular}{p{4cm} p{10cm}}
\toprule
\textbf{Поле} & \textbf{Описание} \\
\midrule
\texttt{user\_id} & Идентификатор пользователя \\
\texttt{courier\_id} & Идентификатор курьера \\
\texttt{order\_id} & Идентификатор заказа \\
\texttt{action} & Действие пользователя или курьера \\
\texttt{time} & Время события \\
\texttt{birth\_date} & Дата рождения курьера \\
\texttt{sex} & Пол курьера \\
\bottomrule
\end{tabular}
\end{table}

В пользовательских действиях встречаются события \texttt{create\_order} и \texttt{cancel\_order}. В действиях курьеров встречаются события \texttt{accept\_order} и \texttt{deliver\_order}. Таблица \texttt{couriers} была использована для дополнительного анализа профиля курьеров: пола, возраста и их связи со временем доставки.

% --------------------------------------------------
\section{Подготовка данных}

Перед анализом данные были загружены в SQLite через DBeaver. Для расчёта времени доставки была создана отдельная таблица \texttt{delivery\_times}, где для каждого заказа сопоставлялись время принятия заказа курьером и время доставки.

После этого была создана очищенная таблица \texttt{delivery\_times\_clean}, из которой были исключены доставки с длительностью 0 минут. Таких записей оказалось всего 3, поэтому они не влияли существенно на общий результат, но были отмечены как возможная особенность логирования данных.

В данных были найдены доставки с длительностью 0 минут. Вероятнее всего, это связано с округлением времени событий до минут или особенностями логирования. Так как таких записей всего 3 на фоне 56 616 доставленных заказов, они не оказывают значимого влияния на выводы.

Для дополнительного анализа профиля курьеров таблица \texttt{delivery\_times\_clean} была объединена с таблицей \texttt{couriers} по полю \texttt{courier\_id}. В анализ пола и возраста были включены только те доставки, для которых удалось сопоставить курьера с его профилем.

После объединения в профильный анализ попало 56 514 доставок и 2 788 активных курьеров. Часть доставок была исключена, так как для них отсутствовала информация о курьере в таблице \texttt{couriers}. Такой подход позволяет корректно сравнивать среднее время доставки по полу и возрастным группам только на записях с полными данными.

Для расчёта возраста дата рождения курьера была преобразована из текстового формата в дату, после чего возраст рассчитывался на дату окончания периода наблюдения -- 8 сентября 2022 года. Затем курьеры были распределены по возрастным группам: \texttt{under 18}, \texttt{18--24}, \texttt{25--34}, \texttt{35--44} и \texttt{45+}.

% --------------------------------------------------
\newpage
\section{Базовые метрики доставки}

Из 59 595 принятых заказов было доставлено 56 616 заказов. Delivery rate составил 95\%.

\begin{table}[H]
\centering
\caption{Базовые метрики доставки}
\begin{tabular}{lr}
\toprule
\textbf{Метрика} & \textbf{Значение} \\
\midrule
Принятые заказы & 59 595 \\
Доставленные заказы & 56 616 \\
Недоставленные заказы & 2 979 \\
Delivery rate & 95\% \\
\bottomrule
\end{tabular}
\end{table}

Таким образом, большая часть заказов успешно доходит до финального этапа доставки. При этом около 5\% заказов не имеют события \texttt{deliver\_order}. Это может быть связано с отменами, техническими сбоями или другими причинами, которые требуют дополнительного анализа.

% --------------------------------------------------
\section{Анализ времени доставки}

Среднее время доставки составило 19,96 минуты. Максимальное время доставки -- 32 минуты.

\begin{table}[H]
\centering
\caption{Статистика времени доставки}
\begin{tabular}{lr}
\toprule
\textbf{Метрика} & \textbf{Значение} \\
\midrule
Количество доставленных заказов & 56 616 \\
Среднее время доставки & 19,96 мин. \\
Максимальное время доставки & 32 мин. \\
\bottomrule
\end{tabular}
\end{table}

\insertfigure
{images/delivery_time_distribution.png}
{Распределение заказов по времени доставки}
{fig:delivery-time-distribution}

Большинство доставленных заказов попадает в интервал от 16 до 30 минут. На эту группу приходится 93,01\% всех доставок. Это говорит о том, что процесс доставки выглядит стабильным, а доля долгих доставок минимальна.

% --------------------------------------------------
\section{Динамика заказов и нагрузки}

\insertfigure
{images/orders_by_day.png}
{Успешные заказы и активные пользователи по дням}
{fig:orders-by-day}

Количество успешных заказов за период наблюдения заметно выросло. В начале периода было 138 успешных заказов, а к 8 сентября их количество достигло 5 474. Пиковое значение наблюдалось 4 сентября -- 6 010 успешных заказов.

Такой рост может говорить об увеличении активности пользователей, расширении пользовательской базы или влиянии внешних факторов, например маркетинговых акций.

\insertfigure
{images/deliveries_by_hour.png}
{Количество доставок по часам}
{fig:deliveries-by-hour}

Наибольшая нагрузка на доставку приходится на вечерние часы. Самые загруженные часы -- с 18:00 до 22:00. Максимальное количество доставок наблюдается в 19:00 -- 3 473 заказа.

\insertfigure
{images/avg_delivery_time_by_hour.png}
{Среднее время доставки по часам}
{fig:avg-delivery-time-by-hour}

Несмотря на вечерний рост количества заказов, среднее время доставки остаётся около 20 минут. Это говорит о том, что курьерская система справляется с повышенной нагрузкой и не допускает резкого ухудшения скорости доставки в часы пик.

% --------------------------------------------------
\section{Анализ эффективности курьеров}

В данных представлено 2 822 курьера. Среднее количество доставок на одного курьера составило 20,06 заказа, максимальное -- 54 заказа.

\begin{table}[H]
\centering
\caption{Общая статистика по курьерам}
\begin{tabular}{lr}
\toprule
\textbf{Метрика} & \textbf{Значение} \\
\midrule
Количество курьеров & 2 822 \\
Среднее количество доставок на курьера & 20,06 \\
Минимальное количество доставок & 1 \\
Максимальное количество доставок & 54 \\
Среднее время доставки по курьерам & 19,95 мин. \\
\bottomrule
\end{tabular}
\end{table}

\begin{figure}[H]
\centering
\includegraphics[width=0.92\textwidth]{images/top_couriers.png}
\caption{Топ курьеров по количеству доставленных заказов}
\label{fig:top-couriers}
\end{figure}

Самые активные курьеры доставили от 49 до 54 заказов. При этом их среднее время доставки находится примерно на уровне общего среднего значения -- около 20 минут. Это значит, что высокая нагрузка на отдельных курьеров не приводит к сильному ухудшению скорости доставки.

Самые быстрые курьеры среди тех, кто выполнил не менее 20 доставок, доставляли заказы в среднем за 18,1--18,6 минуты. Самые медленные -- примерно за 21,3--22,1 минуты. Разница есть, но даже у медленных курьеров доставка остаётся в пределах приемлемого интервала.


% --------------------------------------------------
\section{Анализ профиля курьеров}

Дополнительно был проведён анализ связи профиля курьера со временем доставки. Для этого использовались данные о поле и дате рождения курьеров из таблицы \texttt{couriers}. После очистки и сопоставления с доставками в анализ попало 2 788 активных курьеров и 56 514 доставленных заказов.

Важно отметить, что такой анализ показывает только различия между группами, но не доказывает причинно-следственную связь. На время доставки также могут влиять район, расстояние, время смены, загруженность ресторанов, погодные условия и другие факторы.

\subsection{Пол курьера и время доставки}

\begin{table}[H]
\centering
\caption{Среднее время доставки по полу курьера}
\begin{tabular}{lrrrr}
\toprule
\textbf{Пол} & \textbf{Курьеры} & \textbf{Доставки} & \textbf{Среднее время} & \textbf{Максимум} \\
\midrule
Female & 1 136 & 22 928 & 19,93 мин. & 31 мин. \\
Male & 1 652 & 33 586 & 19,97 мин. & 32 мин. \\
\bottomrule
\end{tabular}
\end{table}

\insertfigure
{images/avg_delivery_time_by_sex.png}
{Среднее время доставки по полу курьера}
{fig:avg-delivery-time-by-sex}

Среднее время доставки у курьеров женского и мужского пола практически не отличается: 19,93 минуты против 19,97 минуты. Разница составляет около 0,04 минуты, то есть примерно 2--3 секунды. По этим данным нельзя говорить о заметной связи между полом курьера и скоростью доставки.

\subsection{Нагрузка на курьеров по полу}

\begin{table}[H]
\centering
\caption{Средняя нагрузка на курьеров по полу}
\begin{tabular}{lrrrr}
\toprule
\textbf{Пол} & \textbf{Курьеры} & \textbf{Заказов на курьера} & \textbf{Минимум} & \textbf{Максимум} \\
\midrule
Female & 1 136 & 20,18 & 1 & 54 \\
Male & 1 652 & 20,33 & 1 & 52 \\
\bottomrule
\end{tabular}
\end{table}

\insertfigure
{images/courier_workload_by_sex.png}
{Среднее количество доставок на курьера по полу}
{fig:courier-workload-by-sex}

Средняя нагрузка на одного курьера также почти одинаковая: 20,18 заказа у female и 20,33 заказа у male. Это говорит о том, что различия в общем количестве доставок в основном связаны с тем, что курьеров male в данных больше, а не с большей индивидуальной нагрузкой.

\subsection{Возраст курьера и время доставки}

\begin{table}[H]
\centering
\caption{Среднее время доставки по возрастным группам}
\begin{tabular}{lrrr}
\toprule
\textbf{Возрастная группа} & \textbf{Курьеры} & \textbf{Доставки} & \textbf{Среднее время} \\
\midrule
under 18 & 6 & 107 & 20,02 мин. \\
18--24 & 283 & 5 522 & 19,98 мин. \\
25--34 & 2 170 & 43 951 & 19,96 мин. \\
35--44 & 300 & 6 307 & 19,94 мин. \\
45+ & 29 & 627 & 19,95 мин. \\
\bottomrule
\end{tabular}
\end{table}

\insertfigure
{images/couriers_by_age_group.png}
{Количество активных курьеров по возрастным группам}
{fig:couriers-by-age-group}

\insertfigure
{images/avg_delivery_time_by_age_group.png}
{Среднее время доставки по возрастным группам}
{fig:avg-delivery-time-by-age-group}

Среднее время доставки по возрастным группам почти не отличается и находится в диапазоне от 19,94 до 20,02 минуты. Самая большая группа курьеров -- 25--34 года: на неё приходится 2 170 активных курьеров и 43 951 доставка. По данным анализа возрастная группа не показывает заметной связи со скоростью доставки.

\subsection{Пол, возраст и время доставки}

\begin{table}[H]
\centering
\caption{Среднее время доставки по полу и возрастной группе}
\begin{tabular}{llrrr}
\toprule
\textbf{Пол} & \textbf{Возраст} & \textbf{Курьеры} & \textbf{Доставки} & \textbf{Среднее время} \\
\midrule
Female & 18--24 & 118 & 2 319 & 19,97 мин. \\
Female & 25--34 & 876 & 17 723 & 19,92 мин. \\
Female & 35--44 & 128 & 2 643 & 19,94 мин. \\
Female & 45+ & 10 & 180 & 19,88 мин. \\
Male & 18--24 & 165 & 3 203 & 19,99 мин. \\
Male & 25--34 & 1 294 & 26 228 & 19,98 мин. \\
Male & 35--44 & 172 & 3 664 & 19,93 мин. \\
Male & 45+ & 19 & 447 & 19,98 мин. \\
\bottomrule
\end{tabular}
\end{table}

\insertfigure
{images/avg_delivery_time_by_sex_age_group.png}
{Среднее время доставки по полу и возрастной группе}
{fig:avg-delivery-time-by-sex-age-group}

При разбиении по полу и возрасту среднее время доставки также остаётся стабильным и находится около 20 минут во всех группах. Небольшие различия между группами не следует интерпретировать как причинную связь. Например, группа female 45+ показывает самое низкое среднее время, но в ней всего 10 курьеров и 180 доставок, поэтому делать сильные выводы по этой группе некорректно.

Итоговый вывод по этому блоку: пол и возрастная группа курьера практически не связаны со средним временем доставки. Во всех группах среднее время находится около 20 минут, а различия составляют всего несколько секунд. Это говорит о стабильности операционного процесса: скорость доставки почти не зависит от демографических характеристик курьера.

% --------------------------------------------------
\newpage
\section{Пользовательские метрики}

Для анализа пользовательского поведения была создана таблица успешных заказов. Успешным считался заказ, который был создан пользователем и не был отменён.

Количество пользователей с успешными заказами составило 21 088. Из них 14 393 пользователя сделали повторный заказ. Repeat order rate составил 68,25\%.

\begin{table}[H]
\centering
\caption{Пользовательские метрики}
\begin{tabular}{lr}
\toprule
\textbf{Метрика} & \textbf{Значение} \\
\midrule
Пользователи с успешными заказами & 21 088 \\
Пользователи с повторными заказами & 14 393 \\
Repeat order rate & 68,25\% \\
\bottomrule
\end{tabular}
\end{table}

Высокий repeat order rate говорит о том, что значительная часть пользователей возвращается в сервис и делает повторные заказы.

\insertfigure
{images/users_by_orders_count.png}
{Распределение пользователей по количеству успешных заказов}
{fig:users-by-orders-count}

Большая часть пользователей делает от 1 до 4 заказов. При этом 6 695 пользователей сделали только один успешный заказ. Это важная зона роста: продукту стоит отдельно работать с пользователями после первого заказа, чтобы увеличивать вероятность повторной покупки.

% --------------------------------------------------
\newpage
\section{Retention-анализ}

Retention-анализ был построен от даты первого успешного заказа пользователя. День 0 -- это день первого успешного заказа, а следующие дни показывают, какая доля пользователей снова совершала успешный заказ.

\begin{table}[H]
\centering
\caption{Retention после первого успешного заказа}
\begin{tabular}{lrr}
\toprule
\textbf{День} & \textbf{Пользователи} & \textbf{Retention} \\
\midrule
Day 0 & 21 088 & 100\% \\
Day 1 & 3 337 & 15,82\% \\
Day 2 & 2 865 & 13,59\% \\
Day 3 & 2 653 & 12,58\% \\
Day 7 & 1 225 & 5,81\% \\
Day 14 & 76 & 0,36\% \\
\bottomrule
\end{tabular}
\end{table}

\insertfigure
{images/retention_curve.png}
{Retention после первого успешного заказа}
{fig:retention-curve}

После первого успешного заказа на следующий день возвращается 15,82\% пользователей. К 7-му дню retention снижается до 5,81\%, а к 14-му дню -- до 0,36\%.

Это говорит о том, что повторные заказы в основном происходят в первые дни после первого заказа. Несмотря на высокий repeat order rate, долгосрочное удержание снижается достаточно быстро.

% --------------------------------------------------
\newpage
\section{Основные выводы}

По итогам анализа можно выделить несколько основных выводов:

\begin{enumerate}
    \item Delivery rate составляет 95\%, что говорит о высокой доле успешно доставленных заказов.
    \item Среднее время доставки составляет около 20 минут.
    \item 93,01\% заказов доставляются в интервале от 16 до 30 минут.
    \item Вечерние часы с 18:00 до 22:00 являются самыми загруженными, но среднее время доставки в этот период почти не растёт.
    \item Курьерская система выглядит стабильной: даже в часы пик скорость доставки остаётся примерно на одном уровне.
    \item Самые активные курьеры доставили от 49 до 54 заказов, при этом их среднее время доставки осталось близким к общему среднему.
    \item Пол курьера практически не связан со средним временем доставки: female -- 19,93 минуты, male -- 19,97 минуты.
    \item Возрастная группа курьера также почти не связана со скоростью доставки: во всех возрастных группах среднее время находится около 20 минут.
    \item Repeat order rate составляет 68,25\%, что говорит о хорошем уровне повторных заказов.
    \item Retention быстро снижается после первой недели: D7 retention составляет 5,81\%.
    \item Пользователи, совершившие только один заказ, являются важной зоной роста для продукта.
\end{enumerate}
% --------------------------------------------------
\newpage
\section{Продуктовые рекомендации}

На основе анализа можно предложить следующие рекомендации:

\begin{enumerate}
    \item Запустить push-уведомления или email-коммуникации после первого заказа.
    \item Протестировать промокод на второй заказ в течение 3--7 дней после первого заказа.
    \item Сделать персональные рекомендации ресторанов на основе первого заказа.
    \item Отдельно анализировать пользователей, которые сделали только один заказ.
    \item Усилить механики возврата после 7-го дня, так как D7 retention составляет только 5,81\%.
    \item Продолжать мониторить вечерние часы пик, так как именно в это время приходится основная нагрузка на доставку.
    \item Изучить причины недоставленных заказов, так как 2 979 заказов не имеют события доставки.
    \item Не делать операционные решения на основе пола или возраста курьеров, так как анализ не показал заметной связи этих признаков со временем доставки.
    \item Для более глубокого анализа скорости доставки добавить данные о районе, расстоянии, погоде, типе транспорта, времени смены и загруженности ресторанов.
\end{enumerate}

% --------------------------------------------------
\newpage
\section{Ограничения анализа}

У анализа есть несколько ограничений:

\begin{itemize}
    \item В датасете нет подробной информации о причинах отмен заказов.
    \item Нет данных о маркетинговых кампаниях, поэтому нельзя точно определить, что влияло на рост заказов по дням.
    \item Нет информации о расстоянии доставки, погоде, типе транспорта и загруженности ресторанов.
    \item Retention рассчитан на ограниченном временном периоде, поэтому долгосрочное удержание стоит анализировать на более длинных данных.
    \item Анализ пола и возраста курьеров показывает только различия между группами, но не доказывает причинно-следственную связь.
    \item Часть доставок была исключена из анализа профиля курьеров, так как для них не удалось сопоставить \texttt{courier\_id} с таблицей \texttt{couriers}.
    \item Некоторые группы курьеров небольшие, например \texttt{under 18} и \texttt{45+}, поэтому выводы по ним нужно интерпретировать осторожно.
\end{itemize}

Эти ограничения не мешают сделать базовые выводы, но они важны для более глубокого продуктового и операционного анализа.

% --------------------------------------------------
\section{Что можно улучшить в дальнейшем}

В дальнейшем проект можно расширить:

\begin{enumerate}
    \item Добавить анализ выручки, среднего чека и ARPU, если в данных есть информация о стоимости заказов.
    \item Построить полноценный когортный retention по неделям.
    \item Сравнить поведение пользователей по сегментам, например по городам или источникам привлечения, если такие поля доступны.
    \item Построить дашборд в Power BI, Tableau или Looker Studio.
    \item Провести анализ отменённых заказов и найти факторы, которые могут быть связаны с отменами.
    \item Добавить данные о расстоянии доставки и проверить, как оно связано со временем выполнения заказа.
    \item Проверить влияние времени суток, района и загруженности ресторанов на скорость доставки.
    \item Провести статистическую проверку различий между группами курьеров, если появится больше признаков и данных.
\end{enumerate}

% --------------------------------------------------
\section{Итог}

Проект показал, что сервис доставки работает достаточно стабильно с точки зрения операционных метрик: большинство заказов доставляется в пределах 30 минут, а среднее время доставки остаётся около 20 минут даже в часы пик.

Дополнительный анализ профиля курьеров показал, что пол и возраст практически не связаны со средним временем доставки. Во всех группах среднее время находится около 20 минут, а различия составляют всего несколько секунд. Это говорит о том, что операционный процесс доставки выглядит достаточно устойчивым и не зависит заметно от демографических характеристик курьера.

С точки зрения продукта, у сервиса хороший repeat order rate, однако retention быстро снижается после первой недели. Поэтому основная зона роста связана не с процессом доставки, а с повторным вовлечением пользователей после первого заказа.

\end{document}
